% GENERATED FILE. Edit canonical lessons, then run npm run generate:books.
% canonical-source-hash: fnv1a64:bd7e3db09d147b44
% canonical-lessons: KA-C40-oota

\chapter{A Meal}
\label{ch:ka-a-meal}
\hlchaptermodality{40}

\begin{chapteropening}
\textbf{By the end of this chapter:} \emph{I can say I am having a meal in Kannada, using the word Chapter 32 used once and never explained, and place it at the hour of two greetings I already know.}

Say \kn{ಊಟ} \kn{ಮಾಡು}, explain ūṭa against everyday tinnu, read the independent vowel \kn{ಊ}, and place the word at the hour of \kn{ಶುಭ} \kn{ಮಧ್ಯಾಹ್ನ} and before \kn{ಶುಭ} \kn{ರಾತ್ರಿ}.
\end{chapteropening}

\section[ūṭa]{\kn{ಊಟ} (ūṭa) — \textquotedblleft{}a meal,\textquotedblright{} the word that closes this chapter's table}

\label{lesson:KA-C40-oota}



\begin{quote}
Chapter 32 taught you \emph{tinnu}, \textquotedblleft{}to eat,\textquotedblright{} and then, in a single line, reached for a different word for a \textbf{full} meal — without ever teaching that word. Fourteen chapters later, here it is.
\end{quote}

\subsection*{You'll want to know: \kn{ಊಟ}}
\begin{quote}
\textbf{\kn{ಊಟ}} — \emph{ūṭa} — \textbf{a meal}
\end{quote}

\begin{quote}
\textbf{\kn{ಊಟ} \kn{ಮಾಡು}.} — \emph{ūṭa māḍu.} — \textquotedblleft{}to have a meal\textquotedblright{} (literally, \textquotedblleft{}meal-do\textquotedblright{})
\end{quote}

\subsection*{The letters in this word}
\emph{(Skip this if you already read Kannada.)} \textbf{\kn{ಊ}} is the \textbf{independent} long-\emph{ū} vowel — the shape a word uses when it \emph{opens} with this sound, the way \textbf{\kn{ಋ}} opened \emph{ṛtu} back in Chapter 14. Every other long-\emph{ū} you have read so far rode on a consonant as the small sign \textbf{\kn{ೂ}}; here, with no consonant in front of it to carry, the vowel stands alone in its full independent form.

\begin{cousinweb}
\textbf{\kn{ಊಟ}} is native Dravidian, a root Tulu keeps as its own everyday word for food, \emph{ūṭa}. Chapter 32 used the phrase \textbf{\kn{ಊಟ} \kn{ಮಾಡು}} to mark a \textbf{full} meal, set against everyday \textbf{\kn{ತಿನ್ನು}} — the same verb Tamil narrowed down to chewing and snacking — but the book left \emph{ūṭa} itself unexplained until today.

A meal keeps close company with greetings you already own. \textbf{\kn{ಶುಭ} \kn{ಮಧ್ಯಾಹ್ನ}} is wished at the very hour of the midday \emph{ūṭa}; \textbf{\kn{ಶುಭ} \kn{ರಾತ್ರಿ}} is wished once the evening one is over. Two greetings that name the hour; one word for what fills it.
\end{cousinweb}

\subsection*{Your turn}
Say these aloud:

\begin{itemize}
\raggedright
  \item \textquotedblleft{}ūṭa māḍu\textquotedblright{} — to have a meal
  \item everyday against full — \textquotedblleft{}tinnuttēne … ūṭa māḍu\textquotedblright{}
  \item the hour and the meal — \textquotedblleft{}śubha madhyāhna … ūṭa\textquotedblright{}
  \item the meal and the goodnight after it — \textquotedblleft{}ūṭa … śubha rātri\textquotedblright{}
  \item this chapter's three drinks, once more — \textquotedblleft{}cahā, kāphi, hālu\textquotedblright{}
\end{itemize}

\begin{tcolorbox}[breakable,colback=teal!4,colframe=teal!35!black,title={Before you move on}]
Say \textquotedblleft{}to have a meal.\textquotedblright{} (\emph{Ūṭa māḍu}.) Which chapter used that exact phrase already, without explaining \emph{ūṭa}? (\textbf{Chapter 32's} \emph{tinnu} lesson.) How does \emph{ūṭa} differ from \emph{tinnu}? (\emph{Tinnu} is \textbf{everyday eating}; \emph{ūṭa} is a \textbf{full meal}.) Name the two greetings that keep company with mealtimes, and which meal each marks. (\textbf{\kn{ಶುಭ} \kn{ಮಧ್ಯಾಹ್ನ}}, the midday meal; \textbf{\kn{ಶುಭ} \kn{ರಾತ್ರಿ}}, after the evening one.)
\end{tcolorbox}
